\chapter{绪论}

\section{研究背景与意义}

\subsection{从能源转型到服务转型：关于汽车角色的重新定义}

过去一百年，汽车的核心价值从未被认真质疑过——把人从 A 点送到 B 点的工具。
内燃机的轰鸣、变速箱的顿挫、油箱的刻度，都在反复强调同一件事：这是一台运动机器。

电动化打破了这个默认逻辑设定，并且这种打破不只是"能源改变"那么简单。
发动机舱让位于电池包，机械联动被线控系统所取代，车辆第一次拥有了足够的车载算力
和稳定的电气架构。这些改变实际上给汽车担负更多的功能扫清了结构上的障碍。
从这个角度看，新能源汽车就是一种可以被编程的移动平台，而不仅仅是更环保的一种车。

目前，"服务化"趋势大多还停留在车机交互和辅助驾驶层面，向物理空间的主动延伸几乎
还没有开始。本文在这个方向上做一次具体的工程探索：以车辆作为地面部署基点，
借助飞行能力把服务范围向车外的物理空间延伸。本文所主要研究的飞行雨伞系统，正是这一思路的具体落脚点。

\subsection{千年未变的发明，与一个真实的需求}

  雨伞是人类历史上最古老的发明之一，其基本形态——一根中轴、数根骨架、一片覆盖——在过去数千年间几乎未曾发生实质性改变。这并非因为人类缺乏创造力，而是在相当长的历史时期内，遮雨这件事情的"体面成本"与普通人的技术可及性之间，始终存在一道难以逾越的鸿沟：你要么淋雨，要么手持一把笨重的伞。
  科技进步带来的一个重要红利，是体面成本的系统性降低。越来越多曾经属于少数人特权的生活品质，正在因为技术的规模化而变得触手可及。人们对出行体验的期待也随之水涨船高——智能汽车已经能够帮助我们在交通过程中抵御恶劣天气，但车门一旦打开，我们便又回到了与坏天气直接对抗的原始状态。这最后几十米的狼狈，是现有技术体系留下的一个真实商业空白。
  传统雨伞并非不够用，而是不够好——它需要你腾出一只手，打乱你的行进节奏，在风雨中与自己和世界笨拙地周旋。飞行雨伞所试图填补的，正是这个空白：让遮雨这件事情从"自己扛着"变成"系统帮你"，让出行的最后一段路，也能保有一份体面与从容。

\subsection{飞行雨伞：不只是一个产品}

  当然，飞行雨伞并不只是一把会飞的伞。
  如果我们把目光稍微放远一些，会发现它更像是一种创新范式的具体例证：将当下成熟的技术积木——嵌入式系统、多旋翼飞行、充气结构、新能源平台——重新组合，对一个已经存在了数千年、从未被认真重新设计过的传统工具进行一次彻底的重构。这种创新不依赖于某项革命性的基础科学突破，而是依赖于对现有技术的重新想象与跨领域整合。
  这或许是当下最值得关注的一类创新路径：我们身边有大量"凑合用了几百年"的工具，它们在过去没有被重新设计，不是因为没有需求，而是因为技术条件尚未就绪。而今天，随着嵌入式计算、新材料、自动控制等领域的系统性成熟，这些工具正在迎来它们各自的"重新发明时刻"。
  本文希望飞行雨伞不仅仅是一个毕业设计的题目，更是对上述创新逻辑的一次小小实践：从一个真实的生活痛点出发，借助现有技术的组合创新，探索一种新的物理服务形态。它是一个开始，关于如何用工程师的语言，认真地回答"生活还能更好吗"这个问题。

\section{国内外研究现状}

\subsection{可变形态飞行机器人研究现状}

可变形态飞行机器人（Morphing/Reconfigurable UAV）是近年来飞行机器人研究中
一个比较活跃的方向，基本思路是让同一架飞行器通过结构变形在不同任务模式
之间切换，而不是为每种任务单独设计一型平台。

在陆空两栖方向上，已有若干工程实现值得关注。
杨帆等人设计了一种倾转翼两栖无人机，把倾转机构和固定翼气动布局结合起来，
实现了旋翼垂直起降与固定翼巡航之间的模式切换，并完成了
地面运动、垂直起降和过渡飞行的完整流程实物验证\cite{yangfan2021}。
李傲翔等人则针对驱动冗余和模式切换不稳定这两个常见问题，
提出了蜗轮自锁倾转机构加单电机双模式复用的变结构方案，
用 SolidWorks 建模和 MATLAB 仿真走通了可行性验证\cite{liaoxiang2025}。
这两项工作在技术路线上有所不同，但都说明单电机复用与自锁机构
在轻量化两栖平台上是走得通的。

在充气/可展开结构方向，吕强等人于2007年在西北工业大学完成了国内较早的充气机翼飞行验证，采用NACA0015翼型的充气机翼实现了连续6分钟的稳定飞行，证明了充气结构在减重与雷达隐身方面的技术可行性\cite{lvqiang2007}。另一家国外研究机构则将充气式再入飞行器技术推进到了实际飞行验证阶段，并进一步拓展了充气展开结构在航空航天领域的应用边界。这些工作共同构成了本文充气伞面结构设计的重要技术参照。

在微型飞行器设计方面，MiniV-Bat双旋翼无人机通过转子盘控制消除非最小相位动态特性，并利用重心偏置实现自稳，在233.7\,g的机体重量下实现了74分钟的悬停续航记录\cite{minivbat2023}，为本系统在紧凑机身约束下追求高效率推进提供了设计参考。

\subsection{多旋翼飞行器建模与控制研究现状}

多旋翼飞行器的动力学建模与控制设计是实现稳定飞行的理论基础。Quan~Quan在其专著中系统阐述了多旋翼飞行器从气动设计、动力学建模到PID控制器综合的完整体系，是国内外多旋翼设计领域引用最广泛的参考文献之一\cite{quanquan2017}。

在抗扰控制方面，Sierra 和 Santos 针对四旋翼在物流任务中质量变化与风扰
共存的问题，提出了基于自适应神经网络估计器的复合控制策略：神经网络在线
估计等效质量扰动与风扰加速度，估计值以前馈方式补偿至 PID 控制器。
他们在风力 6 至 9 级（蒲福风级）、质量三倍突变的极端条件下，
完成了线性、螺旋、圆形、双纽线等多种轨迹的仿真验证\cite{sierra2019}。
本系统同样面临户外复杂风场与开伞瞬间负载突变的场景，
这一工作的控制思路对本文具有直接的参考价值。

\subsection{无人机自主降落与回收平台研究现状}

无人机向移动平台的精确降落，是车载无人机系统实现自主运行绕不开的一个
技术难题，目前已有不少工作从传感器方案和机构设计两个角度切入。

在视觉引导降落方面，Keipour 等人的工作约束条件比较极端：仅使用单目相机、
点激光测距与机载计算机，不依赖 RTK 或运动捕捉，在以 15\,km/h 行驶的
移动车辆上完成了自主降落，22 次室外测试成功 19 次\cite{keipour2022}。
在传感器配置如此受限的情况下能跑通这个速度，说明纯视觉伺服方案
在工程上已经具备一定的实用性。

在降落平台机构设计方面，Galimov 等人做了一个比较系统的综述工作，
把现有无人机降落站的定位机构归纳为主动推杆式、被动漏斗式、
闭合轮廓式和非标准式四类，并从定位精度、优缺点和工程适用性几个
维度做了对比分析\cite{galimov2020}。本文回收平台的结构方案论证
主要参考了这一分类框架。

在车载无人机库结构设计方面，国内已有若干面向工程落地的研究。
谢立威等人用 SolidWorks 完成了一套无人机起降与回收平台的三维设计，
覆盖了升降机构、导引系统、无线充电和太阳能供电模块，
并做了关键部件的参数计算和选型\cite{xieliwei2023}。
李其朋和葛泽南则针对车载无人机库回收系统，提出了水平归中加垂直
剪叉升降的模块化方案，通过 MATLAB 灵敏度分析与多工况有限元仿真
（静力、动力学、随机振动）完成了电缸推力选型与结构优化，
最终用样机实测验证了仿真结果\cite{liqipeng2026}。

\section{核心概念界定}

\subsection{嵌入式机器人系统的定义与范畴}

本文所说的"嵌入式机器人系统"，本质上是一类把计算、传感、驱动和控制都放在本体上的机器人系统。它以嵌入式处理器为核心，不依赖外部服务器，也能自己完成自主或半自主任务。

如果和云端机器人架构相比，这类系统最直接的特点就是：很多关键计算不放在远端，而是留在本地实时完成。这个差别看起来像是架构选择，实际上会直接影响控制效果。尤其对飞行控制这种对延迟很敏感的任务来说，算得够不够快、链路是不是稳定，往往就是系统能不能飞稳的前提\cite{quanquan2017}。

本项目所用到的飞行控制系统以Pixhawk6C开源飞控为主硬件，运行环境为NuttX实时操作系统，主要完成传感器数据融合、姿态估计、控制律解算和电机指令输出等工作。最后选择这套方案，并不是因为它"常用"，而是因为它的确更适合当前的项目，一方面Pixhawk~6C配合NuttX的实时响应能力更适合户外强干扰场景，另一方面整套开源生态已经比较成熟，调试资料、参数经验、问题定位路径都比较完整，实际操作起来会顺很多。

\subsection{本文对"Agent"使用范畴的界定}

智能体（Agent）在人工智能领域一般是指可以感知环境并作出决策的系统
并执行动作的自主实体\cite{russell2020}。大语言模型兴起后，
这个词被使用得越来越广泛，软件系统和物理机器人都可以包含在内。
本文用“Agent”来强调飞行雨伞系统的一个特点，即
它不是被动等待遥控的，而是可以按照用户的指令来执行自主操作
飞行展开与遮雨保护一起完成。这是物理层面的自主性，
和 AI Agent 的认知决策是两回事。

需要说明的是，视觉跟随、用户意图识别这些感知层功能，
以及更完整的自主决策闭环，都不在本文的工程范围之内。
这里做出界定，主要是为了避免读者误解本文的技术范畴。


\section{研究目标与创新点}

\subsection{研究目标}

本文的总体目的就是设计和实现一个车载多形态充气结构
飞行机器人系统（飞行雨伞），从概念设计、理论计算到实施，走完了走完从概念设计、理论计算到
样机制作及实验验证的全部过程，证明充气结构可以用于飞行器
对形态主动重构技术可行性以及工程实现途径进行研究，探究这些技术在个人出行场景中起到怎样的作用。
具体目标分解如下：
\begin{enumerate}
    \item \textbf{形态重构可行性检验}，设计制作充气伞形结构样件
    充气展开之后充气体可以给多旋翼系统提供足够的结构刚度
    排气后能实现紧凑收纳；
    \item \textbf{飞行系统工程实现}，动力系统计算、旋翼布局设计
    与飞控参数整定，使装有充气伞形结构的多旋翼系统可以保持稳定的悬停状态
    \item \textbf{完整工作流程验证}：跑通收纳$\to$充气展开$\to$起飞
    $\to$悬停$\to$降落$\to$收纳的完整闭环，检验各个阶段衔接的
    连贯性与可重复性；
    车载平台方案设计，车载固定结构和
    充气接口与快速部署机构的车载平台方案。
\end{enumerate}

\subsection{主要创新点}

本文的创新不依赖某项新的基础技术，而在于把几块现成的技术
拼在一起，用到一个从来没人认真重新设计过的场景上。
具体体现在以下三点。

\textbf{创新点一：用充气结构驱动飞行器形态主动重构}

把薄壁充气结构引入多旋翼飞行器的形态切换设计，
以充气/排气过程驱动飞行器在收纳态与飞行态之间主动切换。
现有的形态重构方案（机臂折叠、翼型变形等）基本都是刚性机构\cite{liaoxiang2025}，
收纳比有限；充气方案能展开面积量级更大的功能性结构（伞面），
在这类需求下是一条新路子。

\textbf{创新点二：针对伞形充气载体的多旋翼布局专项设计}

充气伞体有几个不寻常的特性：面积大、非刚性、气动阻力显著，
充气过程还会改变转动惯量和重心位置。
现有的多旋翼设计方法论\cite{quanquan2017}并不覆盖这类载体，
本文针对这些特性专项开展了旋翼布局、推重比和重心匹配的设计。

\textbf{创新点三：以车辆为基站的飞行机器人出行服务方案}

把飞行机器人的角色从"独立飞行平台"改成"车载服务系统的
延伸执行器"，以汽车作为充气展开基站与起降平台，
设计了覆盖收纳固定、充气展开、离车服务与归位回收的
完整服务闭环。这一方向与 Galimov 等人综述中指出的
车载无人机站发展趋势一致\cite{galimov2020}。


\section{本文章节安排}

本文共分八章。

\textbf{第一章}（本章）为绪论，介绍研究背景、国内外现状、
核心概念界定及研究目标与创新点。

\textbf{第二章}进行系统需求分析与总体方案设计，分析应用场景
与使用流程，建立关键性能指标体系，论证充气方案与刚性方案的
优劣，提出系统总体架构。

\textbf{第三章}完成充气伞形结构设计，从力学基础出发，
覆盖气囊设计、材料选型、骨架机构设计与充气展开系统实现。

\textbf{第四章}进行动力系统理论计算与选型，建立质量预算
与升力需求模型，运用动量理论与叶素理论完成旋翼气动计算
和动力系统选型。

\textbf{第五章}介绍嵌入式飞行控制系统，建立多旋翼动力学模型，
分析充气结构对控制系统的影响，记录 PID 参数整定过程。

\textbf{第六章}介绍系统集成与车载平台实现，涵盖充气伞体制作工艺、
整机装配过程与车载平台设计。

\textbf{第七章}进行实验验证与分析，对充气结构性能、飞行性能
与完整工作流程分别开展测试，记录迭代优化过程。

\textbf{第八章}为结论与展望，归纳创新点，指出当前局限性，
展望后续工作方向。